\scanpage{9}
contravariant $T$ de $\underline C$ dans $\underline C'$, i.e. d'un foncteur
\[
T:\underline C^{o}\handwritten{\longrightarrow}\underline C'\quad,
\]
(ce cas se ramenant au précédent en remplaçant $\underline C'$ par $ {\underline C'}^{o}$). On trouve
maintenant un homomorphisme canonique, fonctoriel en $M$:
\[
T(M)\handwritten{\longrightarrow}\operatorname{Hom}_A(M,T(A))\quad,
\]
et il faudra expliciter les énoncés duals de 2.6. et 2.7. , en particulier
l'homomorphisme précédent est un isomorphisme sss $T$ est exact à gauche
(i.e. transforme une suite exacte $M''\handwritten{\unsure{-|\to}}M'\handwritten{-}M\handwritten{-}0$ dans $\underline C$ en une suite
exacte $0\handwritten{-}T(M)\handwritten{-}T(M')\handwritten{-}T(M'')$ ): on ob\handwritten{t}ient une \deleted{\unsure{é|e}}équivalence entre la
catégorie des\handwritten{\unsure{/}}foncteurs contravariants exacts à g de $\underline C$ dans $\underline C'$ et la
catégorie $\handwritten{A(}\underline C'\handwritten{)}$. Ces résultats s'étendent \inserted{\handwritten{comme dans 2.4.}} au cas où $\underline C$ est la catégorie de
tous les $A$-modules, sans hypothèse noethérienne sur $A$, pourvu que dans $\underline C'$
les produits quelconques existent.

Nous allons généraliser les considérations précédentes à la catégorie
$\underline C$ de l'exemple \deleted{\unsure{3}}2.4. Soit pour tout $n\handwritten{\geq}0$, $\underline C^n$ la sous-catégorie pleine de
$\underline C$ formée des modules de type fini sur $\underline C$ annulés par $I^{n+1}$, catégorie
équivalente à la catégorie des modules de type fini sur $A_n=A/I^{n+1}$. Alors
$\underline C_n\handwritten{\longrightarrow}\underline C$ est un foncteur exact, et $\underline C$ est réunion filtrante de ses sous-catégories $\underline C_n$. Alors un foncteur \handwritten{additif}
\[
T:\underline C\handwritten{\longrightarrow}\underline C'
\]
est défini par la suite des foncteurs restriction
\[
T_n:\underline C_n\handwritten{\longrightarrow}\underline C'\quad,
\]
et en vertu de ce qui précède, pour chaque $T_n$ on a un homomorphisme
fonctoriel
\[
T_n(A_n)\otimes_{A_n}M\handwritten{\longrightarrow}T_n(M)=T(M)
\]
qu'on peut aussi écrire
\[
\deleted{\unsure{T_n(}}\qquad T(A_n)\otimes_A M\handwritten{\to}T_{\deleted{\unsure{n}}}(M)\quad\handwritten{(n\geq n_{\unsure{o|0}})}\;.
\]
Pour $M$ fixé dans un $\underline C_{n_o}$, on obtient ainsi un système projectif \deleted{de}
\deleted{\unsure{$T_n$}} $(T(A_n)\otimes_A M)_{n\handwritten{\geq}n_o}$ d'objets de $\underline C'$, et un homomorphisme de ce dernier dans
le système projectif constant défini par \deleted{$M$} l'objet $T(M)$. Sous réserve
de l'existence de $\underset{\handwritten{\longleftarrow}}{\lim}T(A_n)\deleted{\unsure{\otimes M}}$, on trouve donc un homomorphisme can
\[
\underset{\handwritten{\longleftarrow}}{\lim}T(A_n)\otimes_A M\handwritten{\longrightarrow}T(M)
\]
d'où par composition un homomorphisme canonique
\[
(\underset{\handwritten{\longleftarrow}}{\lim}T(A_n))\otimes_A M\handwritten{\longrightarrow}T(M)\quad\handwritten{,}
\]
% archive sheet number: 8

